\section{Representasi State}

\section{\textit{Constraint} Permasalahan}

\section{\textit{Succesor Function}}

\section{\textit{Objective Function}}

\section{Formulasi \textit{Initial State}}

\section{Implementasi Hill Climbing}

Seluruh varian \textit{hill climbing} diimplementasikan pada berkas \texttt{src/\allowbreak algorithms/\allowbreak hill\_\allowbreak climbing/\allowbreak algorithm.py} di atas representasi \textit{state}, validator \textit{constraint}, dan fungsi \textit{neighbor} yang sama dengan algoritma lainnya (\texttt{State}, \texttt{validate\_state}, dan \texttt{random\_neighbor}). Dengan demikian seluruh algoritma menerima masukan dan mengembalikan keluaran yang seragam, yaitu objek \texttt{SearchResult} yang memuat \textit{state} awal, \textit{state} akhir, nilai \textit{objective function} akhir, riwayat nilai per iterasi, durasi pencarian, dan metrik khusus algoritma.

\subsection{Pembangkitan \textit{Neighbor} dan Heuristik}
Ruang \textit{neighbor} pada permasalahan ini sangat besar karena koordinat paket berupa bilangan bulat, sehingga seluruh \textit{neighbor} tidak dapat dienumerasi pada setiap iterasi. Sebagai gantinya, setiap \textit{neighbor} dibangkitkan oleh \texttt{random\_neighbor}, yang secara acak memilih salah satu dari tiga operasi (menukar dua paket, memindahkan paket ke koordinat lain, atau merotasi paket $90^{\circ}$) dan hanya mengembalikan \textit{state} yang lolos validasi \textit{constraint}.

Heuristik yang kami gunakan adalah \textit{steepest} berbasis sampel: pada setiap iterasi dibangkitkan \texttt{neighbors\_per\_iteration} (bawaan 30) \textit{neighbor} valid, kemudian \textit{neighbor} dengan nilai terbaik dianggap sebagai \textit{neighbor} terbaik. Konsekuensinya, varian \textit{steepest} dan \textit{sideways} tidak menjamin menemukan \textit{neighbor} terbaik yang sesungguhnya, tetapi peluangnya naik seiring bertambahnya ukuran sampel. Varian \textit{stochastic} memakai tepat satu \textit{neighbor} per iterasi.

\subsection{Parameter dan Kondisi Berhenti}
\begin{table}[h]
    \centering
    \begin{tabular}{@{}p{4.7cm}p{2.3cm}p{6.5cm}@{}}
        \toprule
        Parameter                          & Bawaan            & Keterangan                                                                                           \\
        \midrule
        \texttt{variant}                   & \texttt{steepest} & \texttt{steepest}, \texttt{sideways}, \texttt{stochastic}, atau \texttt{random\_restart}             \\
        \texttt{neighbors\_per\_iteration} & 30                & Jumlah \textit{neighbor} yang disampel per iterasi (tidak dipakai \textit{stochastic})               \\
        \texttt{max\_iterations}           & 1000              & Batas iterasi; menjadi kondisi berhenti pada \textit{stochastic} dan batas pengaman pada varian lain \\
        \texttt{max\_sideways}             & 100               & Batas \textit{sideways move} (varian \textit{sideways})                                              \\
        \texttt{max\_restarts}             & 10                & Jumlah \textit{restart} setelah percobaan pertama (varian \textit{random restart})                   \\
        \texttt{seed}                      & ---               & Seed pembangkit acak agar eksperimen dapat diulang                                                   \\
        \bottomrule
    \end{tabular}
    \caption{Parameter algoritma \textit{hill climbing}}
\end{table}

\begin{itemize}
    \item \textit{Steepest}: berhenti ketika \textit{neighbor} terbaik dari sampel tidak lebih baik daripada \textit{state} saat ini.
    \item \textit{Sideways}: sama seperti \textit{steepest}, tetapi \textit{neighbor} bernilai sama diterima selama jumlah \textit{sideways move} belum mencapai \texttt{max\_sideways}.
    \item \textit{Stochastic}: berhenti setelah \texttt{max\_iterations} iterasi. Iterasi yang \textit{neighbor}-nya ditolak tetap dihitung.
    \item \textit{Random restart}: menjalankan \textit{steepest} dari \textit{state} awal yang diberikan, lalu dari \texttt{max\_restarts} \textit{state} acak baru, sehingga total terdapat \texttt{max\_restarts}\,+\,1 pendakian. Hasil terbaik dari seluruh pendakian dikembalikan.
\end{itemize}

\subsection{Metrik yang Dicatat}
Setiap \textit{run} mengembalikan \texttt{HillClimbingMetrics} yang berisi varian yang dipakai, banyak iterasi hingga berhenti, jumlah dan batas \textit{sideways move}, jumlah dan batas \textit{restart}, serta banyak iterasi pada setiap \textit{restart}. Riwayat nilai \textit{objective function} dicatat pada setiap iterasi (elemen pertama adalah nilai \textit{state} awal) dan digunakan untuk membuat plot nilai \textit{objective function} terhadap banyak iterasi. Durasi pencarian diukur dengan \texttt{perf\_counter}.

\subsection{\textit{Source Code}}
Fungsi \texttt{\_climb} menjalankan satu pendakian untuk varian \textit{steepest}, \textit{sideways}, atau \textit{stochastic}. Fungsi \texttt{hill\_climbing} memvalidasi parameter, memilih varian, menangani pengulangan pada \textit{random restart}, lalu menyusun \texttt{SearchResult}.

\begin{lstlisting}[language=Python, caption={Fungsi bantu dan \texttt{\_climb} pada \texttt{algorithm.py}}]
from math import isfinite
from random import Random
from time import perf_counter

from src.contracts import ObjectiveFunction, resolve_objective
from src.model import State
from src.neighbors import random_initial_state, random_neighbor
from src.results import HillClimbingMetrics, ObjectiveValue, SearchResult

VARIANTS = ("steepest", "sideways", "stochastic", "random_restart")


def _check_int(value, name: str, minimum: int) -> None:
    if type(value) is not int or value < minimum:
        raise ValueError(f"{name} must be an integer >= {minimum}")


def _score(objective: ObjectiveFunction[State], state: State) -> ObjectiveValue:
    score = objective(state)
    if isinstance(score, bool) or not isinstance(score, (int, float)) or not isfinite(score):
        raise ValueError("objective_fn must return finite integer or float values")
    return score


def _climb(
    start: State,
    objective: ObjectiveFunction[State],
    rng: Random,
    *,
    variant: str,
    neighbors_per_iteration: int,
    max_iterations: int,
    max_sideways: int,
) -> tuple[State, ObjectiveValue, list[ObjectiveValue], int, int]:
    """One climb. Returns (state, score, history, iterations, sideways_moves).

    history[0] is the start score; one entry is appended per iteration.
    """
    current, current_score = start, _score(objective, start)
    history = [current_score]
    iterations = sideways = 0

    while iterations < max_iterations:
        if variant == "stochastic":
            candidates = [random_neighbor(current, rng)]
        else:
            candidates = [random_neighbor(current, rng) for _ in range(neighbors_per_iteration)]
        scored = [(_score(objective, c), c) for c in candidates]
        best_score, best = max(scored, key=lambda pair: pair[0])

        if best_score > current_score:
            current, current_score = best, best_score
        elif variant == "sideways" and best_score == current_score and sideways < max_sideways:
            current = best
            sideways += 1
        elif variant == "stochastic":
            iterations += 1  # a rejected random neighbor still costs an iteration
            history.append(current_score)
            continue
        else:
            break  # local optimum (or sideways budget exhausted)
        iterations += 1
        history.append(current_score)

    return current, current_score, history, iterations, sideways
\end{lstlisting}

\begin{lstlisting}[language=Python, caption={Fungsi \texttt{hill\_climbing} pada \texttt{algorithm.py}}]
def hill_climbing(
    initial_state: State,
    *,
    seed: int,
    objective_fn: ObjectiveFunction[State] | None = None,
    variant: str = "steepest",
    neighbors_per_iteration: int = 30,
    max_iterations: int = 1000,
    max_sideways: int = 100,
    max_restarts: int = 10,
    objective_name: str | None = None,
    run_index: int = 1,
) -> SearchResult[State]:
    """Run one hill-climbing variant and return a shared-contract SearchResult."""
    if variant not in VARIANTS:
        raise ValueError(f"variant must be one of {VARIANTS}")
    if type(seed) is not int:
        raise ValueError("seed must be an integer")
    _check_int(neighbors_per_iteration, "neighbors_per_iteration", 1)
    _check_int(max_iterations, "max_iterations", 0)
    _check_int(max_sideways, "max_sideways", 0)
    _check_int(max_restarts, "max_restarts", 0)
    _check_int(run_index, "run_index", 1)

    started_at = perf_counter()
    rng = Random(seed)
    objective = resolve_objective(initial_state, objective_fn)
    climb_kwargs = dict(
        variant="steepest" if variant == "random_restart" else variant,
        neighbors_per_iteration=neighbors_per_iteration,
        max_iterations=max_iterations,
        max_sideways=max_sideways,
    )

    if variant == "random_restart":
        # Restart 0 uses the supplied initial state; each later one a fresh random state.
        best_state, best_score = initial_state, None
        history: list[ObjectiveValue] = []
        per_restart: list[int] = []
        for restart in range(max_restarts + 1):
            start = (
                initial_state
                if restart == 0
                else random_initial_state(initial_state.packages, initial_state.truck, rng)
            )
            state, score, hist, iters, _ = _climb(start, objective, rng, **climb_kwargs)
            history.extend(hist)
            per_restart.append(iters)
            if best_score is None or score > best_score:
                best_state, best_score = state, score
        final_state, final_score = best_state, best_score
        iterations, sideways = sum(per_restart), 0
        restarts = max_restarts
    else:
        final_state, final_score, history, iterations, sideways = _climb(
            initial_state, objective, rng, **climb_kwargs
        )
        per_restart, restarts = [], 0

    parameters: dict[str, int | float | str | bool] = {
        "variant": variant,
        "max_iterations": max_iterations,
    }
    if variant != "stochastic":
        parameters["neighbors_per_iteration"] = neighbors_per_iteration
    if variant == "sideways":
        parameters["max_sideways"] = max_sideways
    if variant == "random_restart":
        parameters["max_restarts"] = max_restarts

    return SearchResult(
        algorithm=f"hill_climbing_{variant}",
        objective_name=objective_name or ("total_value" if objective_fn is None else "custom"),
        initial_state=initial_state,
        final_state=final_state,
        objective_value=final_score,
        objective_history=tuple(history),
        duration_seconds=perf_counter() - started_at,
        metrics=HillClimbingMetrics(
            variant=variant,
            iterations=iterations,
            sideways_moves=sideways,
            maximum_sideways_moves=max_sideways if variant == "sideways" else None,
            restarts=restarts,
            maximum_restarts=max_restarts if variant == "random_restart" else None,
            iterations_per_restart=tuple(per_restart),
        ),
        seed=seed,
        run_index=run_index,
        parameters=parameters,
    )
\end{lstlisting}

\begin{lstlisting}[language=Python, caption={Fungsi pembungkus tiap varian pada \texttt{algorithm.py}}]
def steepest_ascent_hill_climbing(initial_state: State, **kw) -> SearchResult[State]:
    return hill_climbing(initial_state, variant="steepest", **kw)


def sideways_hill_climbing(initial_state: State, **kw) -> SearchResult[State]:
    return hill_climbing(initial_state, variant="sideways", **kw)


def stochastic_hill_climbing(initial_state: State, **kw) -> SearchResult[State]:
    return hill_climbing(initial_state, variant="stochastic", **kw)


def random_restart_hill_climbing(initial_state: State, **kw) -> SearchResult[State]:
    return hill_climbing(initial_state, variant="random_restart", **kw)
\end{lstlisting}

\subsection{Pelaksana Eksperimen dan Visualisasi}
Berkas \texttt{src/\allowbreak experiments/\allowbreak hill\_\allowbreak climbing.py} menjalankan setiap varian sebanyak tiga kali dengan tiga \textit{seed} berbeda. \textit{State} awal setiap \textit{run} dibangkitkan secara acak dari \textit{seed} tersebut. Untuk setiap \textit{run}, program menyimpan berkas \texttt{result.json}, plot nilai \textit{objective function} terhadap iterasi, dan visualisasi tiga dimensi \textit{state} awal dan akhir. Ringkasan lintas \textit{run} (rata-rata dan simpangan baku nilai akhir, durasi, dan iterasi) disimpan pada \texttt{summary.json} dan \texttt{summary.png}. Program dijalankan melalui \texttt{src/\allowbreak experiments/\allowbreak run\_\allowbreak hill\_\allowbreak climbing.py}.

\section{Implementasi Simullated Annealing}

\section{Implementasi Genetic Algorithm}



